\section{SF-E831. Every block must meet enough actual exceptional
support}\label{SF-E831:sf-e831.-every-block-must-meet-enough-actual-exceptional-support}

Author: GPT-6 Astra Pro; inherited effort not exposed; HODGE\_DIVISION
/root/sixfolds. Date: 2026-11-21.

\textbf{Status:} independently whole-reviewed by /root and
/root/g34\_eightfolds; both mathematical reviews PASS on the frozen draft
\texttt{870daf7349d51f72253507005b5048dbf6e3e9e3232d9cbc55175d85d50eed45}.
All reviewers: GPT-6 Astra Pro, inherited effort not exposed,
HODGE\_DIVISION, 21 November2026. Both new dependencies E829/E830 have also
received complete independent whole-proof reviews. The peer supplied a
separate exact exterior/incidence supplement. This is a strict blockwise
strengthening of their common-space bounds. It supplies necessary incidence
and rank conditions on an actual balanced extension; it constructs no
balanced semiregular object.

\subsection{1. Separate support spaces at the actual
vertices}\label{SF-E831:separate-support-spaces-at-the-actual-vertices}

Use the actual coherent extension and its specified coherent finite action
from E830, with actual equivariant blocks

\[
 0\to A\xrightarrow{i}E\xrightarrow{p}B\to0,
 \quad A=A_0\oplus\bigoplus_{a=1}^u L_a^{n_a},\quad
 B=\bigoplus_{b=1}^v M_b^{m_b}.
 \tag{E831.1}
\]

For each actual matrix component e\_ab, choose the subspace Z\_ab from the
full E830 support table. All that is required in the proof is the two
properties proved there: its diagonal products with reverse degree-one
classes lie in Z\_ab, and its entries are killed into degree3 by every
two-vector in Z\_ab perpendicular for the wedge pairing. Put Z\_ab=0 when
e\_ab is zero. Let J be any subset with Ext3(M\_b,A0)=0. Define

\[
 Z_a=\sum_bZ_{ab},\qquad Z_b=\sum_aZ_{ab}\ (b\in J),
 \quad w_v=6-\dim Z_v-
       \dim\operatorname{rad}(q_{\mathrm{wedge}}|_{Z_v}).
 \tag{E831.2}
\]

Here v denotes one selected vertex: every positive line block, or a negative
block in J. There is no identification of Z\_a and Z\_b on different
vertices. Then

\[
 \operatorname{ext}^2_\Gamma(E,E)
      \ge\sum_{a=1}^u w_a+\sum_{b\in J}w_b.
 \tag{E831.3}
\]

This statement depends on the actual extension entries, and not merely on
which Ext1 spaces could be nonzero. A component set to zero does not remove
any diagonal classes in this bound.

\subsection{2. Full proof through the two actual boundary
maps}\label{SF-E831:full-proof-through-the-two-actual-boundary-maps}

For each selected vertex choose W\_v inside Z\_v perpendicular complementary
to its intersection with Z\_v. It has dimension w\_v and intersects Z\_v
trivially. Because the allowed functors act trivially on H2(O\_X), these
choices are automatically invariant. Take

\[
 D_A=\bigoplus_a W_a\operatorname{Id}_{n_a}
           \subset\operatorname{Ext}^2(A,A).
 \tag{E831.4}
\]

The a-th diagonal coefficient of the image of gamma in Ext1(A,B) is sum\_b
e\_ab{[}1{]} gamma\_ba. It lies in Z\_a tensor End(C\^{}\{n\_a\}). Thus that
image has zero intersection with D\_A, separately on each diagonal. The map
Ext2(A,A) to Ext2(A,E) is consequently injective on D\_A.

On d in D\_A, the outgoing boundary from Ext2(A,E) in Hom(-,E) is
i{[}3{]}d{[}1{]}e. Its (a,b) line component is multiplication of the actual
entry e\_ab by the scalar class d\_a in W\_a. Since W\_a lies in Z\_a
perpendicular and Z\_ab is contained in Z\_a, this class is in Z\_ab
perpendicular and kills the entry. The A0 target is zero because d has zero
A0 block. Thus all of D\_A lies in the actual image I of Ext2(E,E) in
Ext2(A,E).

For the selected negative vertices, take

\[
 D_B=\bigoplus_{b\in J}W_b\operatorname{Id}_{m_b}
           \subset\operatorname{Ext}^2(B,B).
 \tag{E831.5}
\]

The boundary to Ext3(B,A) is e{[}2{]}h. For b in J, its line-target
components are killed because W\_b lies in every incident Z\_ab
perpendicular. Its A0 component vanishes by Ext3(M\_b,A0)=0. Hence D\_B
lifts into Ext2(B,E).

Project the preceding boundary from Ext1(A,E) to Ext2(B,B). Its b-diagonal
is the sum of reverse degree-one classes multiplied by the actual e\_ab
entries, with the specified multiplicity-matrix compositions. The line terms
belong to Z\_b tensor End(C\^{}\{m\_b\}). The A0 term is zero by Serre
duality from the same selected vanishing. For each b independently,
normalized matrix trace modulo Z\_b kills this boundary and is the identity
on W\_b times the identity matrix. Therefore the quotient

\[
 Q=\operatorname{Ext}^2(B,E)\big/
 \operatorname{im}\bigl(\operatorname{Ext}^1(A,E)
                   \to\operatorname{Ext}^2(B,E)\bigr)
 \tag{E831.6}
\]

has dimension at least sum\_b in J w\_b and injects into Ext2(E,E). Its
quotient there is I, which already contains D\_A. The two bounds add. Every
boundary, projection and trace is equivariant. The identity multiplicity
matrices are fixed even if their chosen descriptions are projective. Finite
invariants are exact, and the invariant diagonal lifts exist by averaging.
Thus the same argument proves(E831.3) after invariants. No spectral-sequence
degeneration or vanishing of uncomputed higher products is used.

\subsection{3. A necessary coverage condition and a weighted edge
budget}\label{SF-E831:a-necessary-coverage-condition-and-a-weighted-edge-budget}

Suppose the full normalized character is a nonzero multiple of balanced
v\_q, with any rational algebraic B field. The reviewed cap rank and full
supplied invariant CY4 adjunction require invariant Ext2 dimension18 under
weak semiregularity. Put

\[
 N=u+|J|,\qquad
 \ell_v=\dim Z_v+
       \dim\operatorname{rad}(q_{\mathrm{wedge}}|_{Z_v})=6-w_v.
 \tag{E831.7}
\]

The exact necessary inequality is

\[
 \sum_v\ell_v\ge6N-18.
 \tag{E831.8}
\]

In particular a selected vertex with Z\_v=0 contributes six surviving
classes. At most three selected vertices can have no nonzero exceptional
component. Thus at least N-3 selected vertices must actually meet a pair
with Z\_ab nonzero. A merely available but unused Ext1 class does not
qualify. This remains true when other vertices have full rank-three elliptic
components: they cannot erase classes at an unrelated vertex.

For an explicit coarse budget put r\_ab=dim Z\_ab. Since rad(q\textbar Z\_v)
is a subspace of Z\_v,

\[
 \ell_v\le2\dim Z_v\le2\sum_{ab\text{ incident to }v}r_{ab}.
 \tag{E831.9}
\]

Each pair is incident to its positive vertex, and also to its negative
vertex exactly when b belongs to J. Consequently

\[
 \sum_{a,b}(1+\mathbf1_{b\in J})r_{ab}\ge3N-9
 \tag{E831.10}
\]

is another necessary condition for weak balance. The precise vertexwise
bound(E831.8) is stronger and remains available whenever the actual
subspaces can be computed.

In the opposite-slope setting, numerical cross differences are nonzero. The
table in E830 then has r\_ab at most3: a surface index-one component has
r\_ab=1, a positive elliptic component has r\_ab equal to its actual kernel
coefficient rank, and all harmless pairs have r\_ab=0. An active pair meets
at most two selected vertices. If N exceeds3, at least ceil((N-3)/2) active
exceptional pairs are necessary, even allowing rank-three elliptic entries.
More precisely the one- or two-endpoint weight in(E831.10) should be used
when some negative vertices are not in J. If all active exceptional pairs
are surface components, or elliptic components of coefficient rank at most1,
then r\_ab at most1 and the number of pairs must be at least ceil(3(N-3)/2).
These bounds are stated only for N greater than3; for smaller N they give no
additional condition.

For the32 distinct line blocks in the E826 example, if all those negative
blocks remain selected, the necessary coarse budgets would be at least15
active maximal elliptic pairs, or at least44 pairs with dim Z\_ab=1. The
original E826 pair differences are nondegenerate and have no such support.
These integers are consequences of the general inequality, not an assertion
that the proposed exceptional pairs can be realized while preserving its
full character or its selected unit vanishings.

\subsection{4. Correlated surface directions needed to remove every local
class}\label{SF-E831:correlated-surface-directions-needed-to-remove-every-local-class}

There is also a useful exact incidence test when dim Z\_v is at most3. Since
its radical has dimension at most dim Z\_v, w\_v=0 implies dim Z\_v at
least3. Under this dimension bound it is therefore equivalent to dim Z\_v=3
and q\_wedge vanishing identically on Z\_v. Such a three-space is maximal
isotropic in the six-dimensional wedge space.

For completeness its elementary classification is as follows. Every nonzero
element z in it is decomposable, because on a four-dimensional space z wedge
z=0 is equivalent to skew rank at most2. Choose two independent elements a
wedge b and a wedge c: their support planes intersect because their mutual
wedge is zero, and independence makes a,b,c linearly independent. If every
support plane contains a, the three-space equals a wedge V. Otherwise choose
a third support plane not containing a. Its intersection lines with
span(a,b) and span(a,c) are distinct and span it, so it lies in
H=span(a,b,c). Its two-vector has a nonzero b wedge c coefficient, and the
whole three-space equals Lambda2 H. Thus the only possibilities are

\[
 Z_v=a\wedge V\quad(\dim\mathbf Ca=1),
 \qquad\hbox{or}\qquad
 Z_v=\Lambda^2H\quad(\dim H=3).
 \tag{E831.11}
\]

Conversely both displayed spaces are three-dimensional and totally
isotropic, so both give w\_v=0. This is a complete linear-algebra proof; no
unproved general-position assertion is used.

In the surface-only case, if at most three independent active support lines
meet a vertex, removing all six diagonal directions therefore requires their
support planes to be of this correlated type: a common one-dimensional
direction, or containment in a common three-dimensional direction. Two
complementary surface planes do not suffice, since their wedge pairing is
nonzero and they leave w\_v=4. If more than three independent support lines
meet the vertex, use(E831.2); the above three-space classification is not
silently extended to that case. These are cohomological support planes from
actual degenerate line forms, not irrational RM eigenspaces. For elliptic
entry coefficients a support subspace may be complex rather than rational;
no extra rationality is inferred from the classification.

\subsection{5. Scope and verification
boundary}\label{SF-E831:scope-and-verification-boundary}

This result applies only when the actual equivariant filtration and line
blocks are present. E826 and E828 explain when their intrinsic HN/evaluation
arguments provide them; a genuinely nonsplit positive factor is still not a
direct sum by assertion. The selected negative vertices continue to require
the exact Ext3(M\_b,A0) vanishing. No numerical count substitutes for that
condition.

The full proof uses the actual-entry support table of E830 and its explicit
finite-isogeny cup-product calculation from E829. The finite-action,
published BF and supplied S21 statuses are exactly those already stated
there. There is no new external classification premise. The exact script
rm\_vertexwise\_wedge\_certificate.py actually passed nine wedge-space
tests, a separately labelled hypothetical six-vertex incidence test, and
all97 closed-form budgets for N=4 through100. Runtime was0.0258956120015s.
The initial output-serialization failure for exact SymPy integers is
retained; only an exact-integer JSON encoder was added. These finite checks
supplement the infinite proof above and do not assert geometric
realizability of the hypothetical incidence data. No semiregular seed,
period-map rank or generic extension-product rank is claimed.
